%% kapa: exam=1 round=20 year=2009 period=2 form=A subject=accounting subject_ko=회계학 seq=24 number=24
%% source: Q-Net 감정평가사 기출문제 (cst00309) artlSeq=5024178 file=extracted/1008/5024178/2009년 제20회 감정평가사 1차시험 기출문제/2교시.hwp
%% conversion: kordoc 4.24.51 -> tools/gen1.py (auto); review=proofread-sonnet; flags=none
\begin{question}{24}
20×8년 말에 (주)한국이 소유하고 있는 재고자산과 창고건물이 화재발생으로 인하여 전소되었다. 화재와 관련된 창고건물과 재고자산은 모두 보험에 가입되어 있었다. 창고건물은 20×6년도 초에 구입(취득원가 ￦90,000, 내용연수 10년, 잔존가치는 취득원가의 10\%, 감가상각방법은 정액법)하였으며, 20×8년도 기말재고자산의 장부금액은 ￦15,000이었다. 보험회사는 화재로 인한 손실을 창고건물 ￦70,000, 재고자산 ￦17,000으로 인정하고, 이 금액을 20×8년 말에 보험금으로 지급하였다. 화재와 보험금 수령이 20×8년도 당기순이익과 자본잉여금에 미치는 영향은?

{\small\centering
\begin{tabular}{@{}ccc@{}}
 & 당기순이익 & 자본잉여금 \\
\cline{2-2}\cline{3-3}
① & ￦2,000 증가 & ￦4,300 증가 \\
② & ￦4,300 증가 & ￦2,000 증가 \\
③ & ￦6,300 증가 & - \\
④ & ￦13,300 증가 & - \\
⑤ & - & ￦6,300 증가 \\
\end{tabular}\par}
\end{question}
\answer{3}
